%% conclusion.tex

\section{Conclusion}
\label{sec:conclusion}

\paragraph{Summary}
This paper introduced \rci{}, a previously overlooked class of attacks against self-healing LLM agent systems. We showed that the recovery pipeline---designed to improve robustness by re-injecting diagnostic context into subsequent iterations---constitutes a \emph{trust-escalation trampoline}: adversarial content that is correctly quarantined as untrusted on the forward execution path is re-laundered as system-trusted on the reverse recovery path, because the recovery machinery was never designed to receive untrusted input. We employed a two-stage methodology: a white-box testbed (\system{}) for mechanism attribution, and a black-box production framework (LangGraph) for external validity.

\paragraph{Mechanism attribution (white-box testbed)}
Because production frameworks do not expose sufficient internal hooks for fine-grained mechanism attribution, we constructed \system{} as a white-box testbed with instrumentation at every trust-boundary crossing. On the testbed, we isolated two attack mechanisms: (V1)~trust escalation via \code{ReflectionInjectionRecovery} and (V2)~privilege escalation via \code{TopologyMutationStrategy}. We emphasize that V2 is not a novel primitive---it is CWE-862 (Missing Authorization), a decades-old vulnerability class---but its \emph{systematic appearance} in self-healing recovery paths is an empirical finding: the recovery pipeline is assumed benign and therefore bypasses the permission checks enforced on the forward path, creating a confused-deputy pattern. In baseline configuration, V2 achieves 100\% ASR and V1 57\% ASR against real LLMs. The testbed enabled controlled ablation of the proposed defense: \trustorigin{} Tagging + \reauthgate{} reduces V1 to 0.07 (Fisher $p = 7.6 \times 10^{-15}$) and V2 to 0.00 (deterministic), with 0.0\% FAR and sub-millisecond overhead. The testbed's role is to establish the \emph{mechanism}---isolating content-level vs.\ structural-level contributions---not to serve as evidence that production systems are vulnerable.

\paragraph{External validity (LangGraph + wild-case study)}
The testbed establishes the mechanism; LangGraph establishes that it is not an artifact of the authors' codebase. A forward-path experiment on \code{create\_react\_agent} using entirely library-default components achieves 0.88 V1 pooled ASR, confirming that production frameworks lack provenance tagging---though this is the classic indirect injection vector, not RCI per se. A \emph{real agentic-loop} experiment (where the LLM actually executes the canary tool rather than describing intent) confirms ecological validity: baseline ASR rises to 0.96 ($[0.80, 0.99]$), validating that the proxy measurement is conservative. A recovery-path experiment on \code{StateGraph} (author-constructed recovery node) achieves 0.76 V1 pooled ASR, reduced to 0.04 by the untrusted-framing defense, demonstrating the RCI mechanism is realizable on a third-party API; V2 achieves 0.54, reduced to 0.00 by the \reauthgate{} analog. A \emph{critical experiment} (\S\ref{sec:forward-defense-bypass}) establishes RCI as an independent attack surface: with three forward-path defenses fully enabled (input sanitizer + LLM injection detector + tool-permission check), the attack still achieves 0.68 ASR ($p=0.660$ vs.\ undefended), because the defenses inspect user input while the payload enters via the tool's failure. A source-level wild-case study on four production agent codebases (Reflexion, MetaGPT, Aider, OpenHands; $\sim$105k stars) found three of four ship retry logic with the V1 trampoline signature. For Reflexion, we additionally confirmed the vulnerability \emph{end-to-end} (\S\ref{sec:reflexion-e2e}): an exploit against Reflexion's actual \code{generic\_generate\_func\_impl} retry path achieves 0.52 ASR ($p = 3.5 \times 10^{-10}$ vs.\ 0\% FPR control), with two variants at 100\%---using Reflexion's verbatim prompt construction, not an author-constructed recovery node. A stronger-model evaluation on gpt-oss-120b reveals content-level V1 defense degrades (0.07 $\rightarrow$ 0.50) while structural V2 defense remains at 0.00. A supplementary evaluation on a closed-weight model (accessed via the official API) finds defended ASR of 0.36 with $p=0.027$ (\emph{significant}, $N=50$): this \emph{confirms} that the content-level V1 defense degrades on this model---the defense reduces ASR (0.59 $\rightarrow$ 0.36) but the residual is far higher than on DeepSeek/Qwen (0.07). The V2 defense is LLM-independent by construction.

\paragraph{Where the burden of proof rests.}
We are explicit about the evidentiary weight of each component, because the testbed is the authors' own prototype and a reviewer may otherwise overweight its numbers. The \emph{burden of proof} for the claim that RCI is a real attack surface rests on the two pieces of evidence that use \emph{no author-constructed code}: (i)~the LangGraph forward-path experiment on entirely library-default \code{create\_react\_agent} components (0.88 ASR), and (ii)~the confirmed end-to-end exploit on Reflexion's actual \code{generic\_generate\_func\_impl} retry path at commit \texttt{218cf0e} (0.52 ASR, $p = 3.5 \times 10^{-10}$ vs.\ 0\% FPR control). The \system{} testbed's role is strictly \emph{mechanism attribution}---isolating content-level (V1) vs.\ structural-level (V2) contributions via controlled ablation---not to serve as evidence that production systems are vulnerable. A reader who is skeptical of the testbed's numbers can disregard them entirely and the paper's central claim still stands on LangGraph + Reflexion alone.

\paragraph{Limitations and future work}
Our trust-tagging mechanism is currently text-oriented: provenance is assigned at the capture site by inspecting the failing tool, and the \code{TaggedContent} record carries a textual payload. In multimodal recovery scenarios---e.g., an agent that fails while parsing a malicious image, an adversarial audio snippet, or a poisoned embedded chart---the untrusted signal may be embedded in a non-textual modality that is later transcribed, OCR'd, or described by a vision model before reaching the recovery pipeline. Once transcribed, the original modality boundary is lost, and a purely text-bound \code{origin} tag may be insufficient to track provenance across the modality conversion.

\textbf{Cross-modal provenance preservation.} Concretely, consider a recovery path where a vision-language model transcribes an adversarial image into text before the diagnostic LLM consumes it. The transcription step destroys the modality boundary: the resulting text is indistinguishable from a tool-generated diagnostic, so the \code{TOOL\_OUTPUT\_EXTERNAL} tag cannot be propagated by the current capture-site inspection (which matches on the failing tool's name, not the content's modality history). We sketch three directions for extending \trustorigin{} to multimodal settings. (i)~\emph{Taint tracking through transcription}: instrument the vision model's output so that any text derived from an external image carries a \code{DERIVED\_FROM\_EXTERNAL\_MODALITY} tag, analogous to dynamic taint tracking in language runtimes; the recovery pipeline then inherits the tag regardless of which tool performed the transcription. (ii)~\emph{Modality-aware capture-site inspection}: extend the capture site to inspect not just the failing tool's name but its declared I/O modalities (e.g., a \code{vision\_parser} tool that declares \code{input\_modality=image, output\_modality=text} would be auto-tagged \code{EXTERNAL} at registration time). (iii)~\emph{Cryptographic provenance}: attach a signed provenance certificate to each piece of content at the boundary crossing, so that downstream recovery nodes can verify the content's origin even after arbitrary transformations. Each direction has trade-offs: taint tracking requires runtime instrumentation, modality-aware inspection requires tool cooperation, and cryptographic provenance adds overhead. A complete multimodal provenance system likely combines all three, and is an open problem we leave to future work.

A second limitation is that our gate enforces a \emph{static} privilege invariant: it checks that the proposed role is in the legitimate downgrade set, but does not inspect \emph{why} the recovery strategy proposed it. A sufficiently sophisticated payload could induce a strategy to propose an action that is individually permissible yet collectively harmful (e.g., a sequence of individually-legitimate downgrades that compose into an unintended capability). As future work, we propose \emph{intent auditing via inner-monologue inspection}: instrumenting the diagnostic LLM's chain-of-thought to verify that the proposed recovery action is justified by the failure's actual cause, rather than by injected reasoning. This would add a semantic, model-level defense layer atop the syntactic, OS-level invariants introduced here---complementary rather than a replacement, and aligned with the broader trend of using LLMs to audit LLM behavior.

A third limitation concerns the scope of the empirical evaluation: our real-LLM trials span two model families (DeepSeek and Qwen) at temperature~0.7, with a stronger open-weights model (gpt-oss-120b) as a third data point. Closed-weight frontier models (GPT-4o, Claude, Gemini), lower-temperature deployments, or agentic frameworks with different tool-calling conventions may exhibit different baseline ASR distributions---in particular, the 0.07 residual ASR on V1 reflects the probabilistic nature of content-based provenance signaling and may shift with model scale. Extending the evaluation to a broader model zoo---including closed-weight frontier systems---and to \emph{fully adaptive multi-round adversaries} that observe defense rejections and iteratively rewrite payloads (our defense-aware experiment in \S\ref{sec:adaptive} is single-round and thus a lower bound) is immediate future work.
